\scanpage{18}
\mspage{8}
fibres géom. de $(\widehat X,\underline I)$ forment les fibres d'un schéma en maquettes \inserted{(ou en MD-graphes)} sur $\unsure{\mathcal S|S}$
\[
(\deleted{\unsure{\underline G|\underline S}}\underline S,\underline{\widetilde A},\sigma_{\underline{\widetilde A}},\underline I,\ \underline{\widetilde A}\to\underline S,\ \underline I\to\underline S,\ \inserted{\underline S\xrightarrow{\underline g}\mathbb N_{\mathcal S}})
\]
(système de rev. finis étales \inserted{de $\mathcal S$} et de morphismes entre ceux-ci), loc. isomorphes à la maquette $G$ donnée. \deleted{\unsure{Le}} On appelle \inserted{$G$-}épinglage de $(\widehat X,\underline I)$ tout isomorphisme entre $G_{\mathcal S}$ et $\underline G(\widehat X,\underline I)$. Si \deleted{$\unsure{\Gamma}$}
\[
(18)\qquad\Gamma=\operatorname{Aut}G
\]
(groupe fini), les $G$-épinglages de $(\widehat X,\underline I)$ s'identifient aux sections d'un certain $\Gamma_{\mathcal S}$-torseur, \deleted{\unsure{qui est le|qui sera le}}\inserted{appelé} \uline{torseur des $G$-épinglages} de $(\widehat X,\underline I)$.

Considérons, sur une base $\mathcal S$ fixée, la catégorie \inserted{(pour les iso)} des courbes standart $G$-épinglées. Pour $\forall\ \alpha\in S$ (ens. des sommets de la maquette) soit $\widehat I_\alpha=I_\alpha\amalg\widetilde A_\alpha$ l'ensemble (de cardinal $\widehat\nu_\alpha=\text{poids total de }\alpha$) somme
\mspage{17}
